\documentclass[11pt,a4paper]{article}

\usepackage[T1]{fontenc}
\usepackage[utf8]{inputenc}
\usepackage[french]{babel}
\usepackage{lmodern}
\usepackage{geometry}
\usepackage{booktabs}
\usepackage{amsmath,amssymb}
\usepackage{enumitem}
\usepackage{xcolor}
\usepackage{hyperref}
\geometry{margin=2.3cm}
\hypersetup{colorlinks=true,linkcolor=blue,urlcolor=blue}

\title{Analyse discriminante prédictive\\
       \large Conception et implémentation du package R \texttt{ldapkg}}
\author{Abdoul-Faisal Ouedraogo}
\date{M2 SISE -- Programmation R -- 2026--2027}

\begin{document}
\maketitle
\tableofcontents
\newpage

\section{Objectif et positionnement}

Le projet consiste à fournir une analyse discriminante linéaire prédictive
(LDA) utilisable comme un package R installable depuis GitHub. L'objectif
est double : produire une décision de classe pour de nouveaux individus et
fournir des sorties interprétables sur la séparation des classes, les
variables et les axes discriminants.

L'interface publique est une classe R6 inspirée de scikit-learn :
\texttt{LDAClassifier\$new()}, \texttt{\$fit(X, y)}, \texttt{\$predict(X)}
et \texttt{\$predict\_proba(X)}. Les méthodes \texttt{\$summary()},
\texttt{\$plot()}, \texttt{\$cv()} et \texttt{\$transform()} complètent
l'utilisation analytique et prédictive.

\section{Modèle statistique}

On note $K$ le nombre de classes, $n_k$ l'effectif de la classe $k$,
$m_k$ son vecteur moyen et $m$ la moyenne globale. Les matrices de sommes
des carrés et produits croisés intra-classes et inter-classes sont
\[
 E = \sum_{k=1}^K\sum_{i\in k}(x_i-m_k)(x_i-m_k)^\top,\qquad
 B = \sum_{k=1}^K n_k(m_k-m)(m_k-m)^\top.
\]
La matrice totale est $T=E+B$ et la covariance poolée est
$W=E/(n-K)$.

Pour une observation $x$, la fonction de classification de la classe $k$ est
\[
 \delta_k(x)=x^\top W^{-1}m_k
 -\frac12m_k^\top W^{-1}m_k+\log(\pi_k),
\]
où $\pi_k$ est la probabilité a priori. La classe prédite maximise
\(\delta_k(x)\). Les probabilités a posteriori sont obtenues par une
normalisation softmax stable des scores.

\section{Architecture du package}

\begin{center}
\begin{tabular}{ll}
\toprule
\textbf{Module} & \textbf{Responsabilité} \\
\midrule
\texttt{preprocess.R} & validation de l'encodage et variables catégorielles \\
\texttt{core.R} & statistiques suffisantes, Cholesky, Wilks, axes \\
\texttt{selection.R} & sélection forward, backward et stepwise \\
\texttt{LDAClassifier.R} & classe R6, apprentissage et prédiction \\
\texttt{methods.R} & méthodes S3 \texttt{predict}, \texttt{summary}, \texttt{plot} \\
\texttt{examples.R} & chargement des jeux de données exemples \\
\texttt{run\_app.R} & lancement de l'application Shiny \\
\bottomrule
\end{tabular}
\end{center}

Le package est limité à des dépendances légères pour le coeur
(\texttt{R6}, \texttt{stats}, \texttt{graphics}, \texttt{grDevices} et
\texttt{utils}). Les dépendances de l'application, comme \texttt{shiny},
\texttt{DT} et \texttt{readxl}, sont optionnelles.

\section{Algorithmes et choix d'implémentation}

\subsection{Calcul efficace}

Les moyennes par classe sont obtenues avec \texttt{rowsum()} et les matrices
avec \texttt{crossprod()}. Le calcul ne construit donc pas une liste de
produits individu par individu. Les résolutions matricielles utilisent des
factorisations de Cholesky plutôt qu'une inversion générale.

Les colonnes constantes et les colonnes linéairement dépendantes sont
détectées avant la factorisation. Elles sont retirées avec un avertissement
explicite. Les variables sont ensuite standardisées sur l'échelle intra-classe,
ce qui améliore le conditionnement sans modifier les prédictions de la LDA.

\subsection{Variables catégorielles}

Une variable factorielle ou caractère est codée par $q-1$ indicatrices, la
première modalité étant la référence. La spécification (noms et niveaux) est
conservée dans le modèle. À la prédiction, les nouvelles données sont donc
encodées dans le même ordre. Un niveau inconnu ne provoque pas de plantage :
la ligne concernée est signalée et reçoit une sortie \texttt{NA}.

\subsection{Sélection de variables}

La sélection travaille sur les matrices $E$ et $T$ déjà calculées. Pour une
variable candidate, le complément de Schur fournit la contribution
conditionnelle. Le test repose sur l'approximation $F$ de Rao associée à la
lambda de Wilks. Les variables catégorielles sont toujours traitées comme
un bloc de colonnes indicatrices.

\begin{enumerate}[label=\arabic*.]
\item Initialiser l'ensemble sélectionné selon la stratégie choisie.
\item Calculer la lambda et la p-valeur de chaque candidate.
\item Ajouter la meilleure candidate si elle respecte \texttt{alpha\_enter}.
\item Pour la stratégie stepwise, retirer la variable la moins utile si sa
      p-valeur dépasse \texttt{alpha\_stay}.
\item Arrêter lorsqu'aucune modification n'est possible.
\end{enumerate}

\subsection{Validation}

\texttt{\$cv()} réalise une validation croisée stratifiée. Le nombre de plis
est automatiquement limité à l'effectif de la plus petite classe afin que
chaque apprentissage conserve toutes les classes. La sélection de variables
est recalculée dans chaque pli, ce qui évite une fuite d'information.

\section{Sorties interprétables}

Le modèle expose les coefficients et constantes des fonctions de
classification, les probabilités a priori, les effectifs et moyennes de
classe, ainsi que les variables et colonnes retenues. La lambda de Wilks
\(\det(E)/\det(T)\), le test $F$ de Rao, les lambdas partielles, les
p-valeurs et les tolérances permettent d'étudier la contribution des
variables.

Les axes canoniques sont calculés par décomposition singulière et normalisés
pour avoir une variance intra-classe unité. Le modèle expose les valeurs
propres, les corrélations canoniques, les coordonnées des individus et les
centroïdes de classe. Les performances comprennent la matrice de confusion,
l'accuracy, le rappel, la précision, le score F1 et la validation croisée.

\section{Application Shiny}

L'application permet de charger un fichier tabulé, CSV ou Excel, ou un jeu de
données exemple. L'utilisateur choisit la cible, les variables actives et les
variables illustratives, puis règle les a priori, la sélection, la taille du
jeu test et le nombre de plis. Les résultats sont présentés sous forme de
tableaux, graphiques et fonctions de classification.

Le déploiement est disponible pour un individu saisi dans un formulaire et
pour un fichier d'individus supplémentaires. Les probabilités par classe
peuvent être téléchargées avec la classe prédite. Les erreurs de lecture, les
niveaux inconnus et les valeurs manquantes sont affichés sous forme de
notifications plutôt que de faire disparaître silencieusement le résultat.

\section{Vérification et reproductibilité}

Les tests de \texttt{testthat} comparent les probabilités et les axes à
\texttt{MASS::lda()}, la lambda à \texttt{stats::manova()}, et couvrent les
variables catégorielles, les valeurs manquantes, les colonnes redondantes,
la sélection, la validation croisée et les prédictions inconnues.

Les commandes de vérification sont :
\begin{verbatim}
devtools::document()
devtools::test()
devtools::check()
devtools::build()
\end{verbatim}

Le workflow GitHub Actions exécute automatiquement \texttt{R CMD check} sur
Linux, Windows et macOS. Le tutoriel anglais dans \texttt{tutorial/} est
reproductible avec les jeux de données fournis dans
\texttt{inst/extdata/}.

\section{Limites et pistes d'amélioration}

La méthode suppose une covariance commune aux classes et reste sensible aux
échantillons très petits. Pour des données de très grande dimension, une
régularisation de la covariance ou une réduction de dimension préalable
serait pertinente. Une extension future pourrait proposer un paramètre de
shrinkage, des courbes ROC et une calibration des probabilités, tout en
conservant l'interface actuelle.

\section{Références}

\begin{itemize}
\item Fisher, R. A. (1936). The use of multiple measurements in taxonomic
      problems. \emph{Annals of Eugenics}, 7, 179--188.
\item Rao, C. R. (1951). An asymptotic expansion of the distribution of
      Wilks' criterion. \emph{Bulletin of the International Statistical
      Institute}, 33, 177--180.
\item Venables, W. N. and Ripley, B. D. (2002).
      \emph{Modern Applied Statistics with S}. Springer.
\item R Core Team (2026). \emph{R: A Language and Environment for Statistical
      Computing}. R Foundation for Statistical Computing.
\end{itemize}

\end{document}
